\section*{結言}
今回の8件では、新たな大型発表の数よりも、既存技術を実際に使うための条件が目立った。Nanbeige4.3-33Bの自己報告ベンチマークはエージェント用途への期待を集めたが、同時に必要メモリや長文処理時の負荷も論点となった。オープンウェイトの価値は、モデルの公開だけでなく、手元の環境で動くかどうかによって左右される。

Claude Duetの継続的な話題は、指示を待つだけでなく日常の作業を先回りして支援するエージェントへの関心を示している。一方、IPOやモデルの提供時期に関する推測は、公式発表と区別して扱う必要がある。ゲーム制作への生成AI活用や広告の透明性を巡る話題からは、技術以外の社会的な論点も見えてきた。

Grok自身が今回のカバレッジをPartialと評価しており、いくつかの投稿時刻は観測窓の外にある。今後は独立した性能測定、製品の提供条件、規制の具体的な施行情報が、今回の議論を評価する材料となる。
